\section{Further Topics in Group Theory}

\subsection{$p$-Groups, Nilpotent Groups, and Solvable Groups}

\begin{exercise}{6.1.1}
    Prove that $Z_i(G)$ is a characteristic subgroup of $G$ for all $i$.
\end{exercise}

\begin{exercise}{6.1.2}
    Prove Parts 2 and 4 of Theorem 1 for $G$ a finite nilpotent group, not necessarily a $p$-group.
\end{exercise}

\begin{exercise}{6.1.3}
    If $G$ is finite, prove that $G$ is nilpotent if and only if it has a normal subgroup of each order dividing $|G|$, and is cyclic if and only if it has a unique subgroup of each order dividing $|G|$.
\end{exercise}

\begin{exercise}{6.1.4}
    Prove that a maximal subgroup of a finite nilpotent group has prime index.
\end{exercise}

\begin{exercise}{6.1.5}
    Prove Parts 2 and 4 of Theorem 1 for $G$ an infinite nilpotent group.
\end{exercise}

\begin{exercise}{6.1.6}
    Show that if $G/Z(G)$ is nilpotent then $G$ is nilpotent.
\end{exercise}

\begin{exercise}{6.1.7}
    Prove that subgroups and quotient groups of nilpotent groups are nilpotent (your proof should work for infinite groups).
    Give an explicit example of a group $G$ which possesses a normal subgroup $H$ such that both $H$ and $G/H$ are nilpotent but $G$ is not nilpotent.
\end{exercise}

\begin{exercise}{6.1.8}
    Prove that if $p$ is a prime and $P$ is a non-Abelian group of order $p^3$ then $|Z(P)| = p$ and $P/Z(P) \cong Z_p \times Z_p$.
\end{exercise}

\begin{exercise}{6.1.9}
    Prove that a finite group $G$ is nilpotent if and only if whenever $a, b \in G$ with $(|a|, |b|) = 1$ then $ab = ba$.
    [Use Part 4 of Theorem 3.]
\end{exercise}

\begin{exercise}{6.1.10}
    Prove that $D_{2n}$ is nilpotent if and only if $n$ is a power of $2$.
    [Use Exercise 9.]
\end{exercise}

\begin{exercise}{6.1.11}
    Give another proof of Proposition 5 under the additional assumption that $G$ is Abelian by invoking the Fundamental Theorem of Finite Abelian Groups.
\end{exercise}

\begin{exercise}{6.1.12}
    Find the upper and lower central series for $A_4$ and $S_4$.
\end{exercise}

\begin{exercise}{6.1.13}
    Find the upper and lower central series for $A_n$ and $S_n$, $n \geq 5$.
\end{exercise}

\begin{exercise}{6.1.14}
    Prove that $G^i$ is a characteristic subgroup of $G$ for all $i$.
\end{exercise}

\begin{exercise}{6.1.15}
    Prove that $Z_i(D_{2^n}) = D_{2^n}^{n-1-i}$.
\end{exercise}

\begin{exercise}{6.1.16}
    Prove that $\qbb$ has no maximal subgroups.
    [Recall Exercise 21, Section 3.2.]
\end{exercise}

\begin{exercise}{6.1.17}
    Prove that $G^{(i)}$ is a characteristic subgroup of $G$ for all $i$.
\end{exercise}

\begin{exercise}{6.1.18}
    Show that if $G'/G''$ and $G''/G'''$ are both cyclic then $G'' = 1$.
    [You may assume $G''' = 1$. Then $G/G''$ acts by conjugation on the cyclic group $G''$.]
\end{exercise}

\begin{exercise}{6.1.19}
    Show that there is no group whose commutator subgroup is isomorphic to $S_4$.
    [Use the preceding exercise.]
\end{exercise}

\begin{exercise}{6.1.20}
    Let $p$ be a prime, let $P$ be a $p$-subgroup of the finite group $G$, let $N$ be a normal subgroup of $G$ whose order is relatively prime to $p$ and let $\bar{G} = G/N$.
    Prove the following:
    \begin{subexercises}
        \item $N_{\bar{G}}(\bar{P}) = \bar{N_G(P)}$ [Use Frattini's Argument.]
        \item $C_{\bar{G}}(\bar{P}) = \bar{C_G(P)}$. [Use part (a).]
    \end{subexercises}
\end{exercise}

\begin{exercise}{6.1.21}
    Prove that $\Phi(G)$ is a characteristic subgroup of $G$.
\end{exercise}

\begin{exercise}{6.1.22}
    Prove that if $N \nsub G$ then $\Phi(N) \leq \Phi(G)$.
    Give an explicit example where this containment does not hold if $N$ is not normal in $G$.
\end{exercise}

\begin{exercise}{6.1.23}
    Compute $\Phi(S_3)$, $\Phi(A_4)$, $\Phi(S_4)$, $\Phi(A_5)$ and $\Phi(S_5)$.
\end{exercise}

\begin{exercise}{6.1.24}
    Say an element $x$ of $G$ is a \emph{non-generator} if for every proper subgroup $H$ of $G$, $\langle x, H \rangle$ is also a proper subgroup of $G$.
    Prove that $\Phi(G)$ is the set of non-generators of $G$ (here $|G| > 1$).
\end{exercise}

\begin{exercise}{6.1.25}
    Let $G$ be a finite group.
    Prove that $\Phi(G)$ is nilpotent.
    [Use Frattini's Argument to prove that every Sylow subgroup of $\Phi(G)$ is normal in $G$.]
\end{exercise}

\begin{exercise}{6.1.26}
    Let $p$ be a prime, let $P$ be a finite $p$-group and let $\bar{P} = P/\Phi(P)$.
    \begin{subexercises}
        \item Prove that $\bar{P}$ is an elementary Abelian $p$-group.
        [Show that $P' \leq \Phi(P)$ and that $x^p \in \Phi(P)$ for all $x \in P$.]
        \item Prove that if $N$ is any normal subgroup of $P$ such that $P/N$ is elementary Abelian then $\Phi(P) \leq N$.
        State this (universal) property in terms of homomorphisms and commutative diagrams.
        \item Let $\bar{P}$ be elementary Abelian of order $p^r$ (by (a)).
        Deduce from Exercise 24 that if $\bar{x_1}, \bar{x_2}, \ldots, \bar{x_r}$ are any basis for the $r$-dimensional vector space $\bar{P}$ over $\fbb_p$ and if $x_i$ is any element of the coset $\bar{x_i}$, then $P = \langle x_1, x_2, \ldots, x_r \rangle$.
        Show conversely that if $y_1, y_2, \ldots, y_s$ is any set of generators for $P$, then $s \geq r$ (you may assume that every minimal generating set for an $r$-dimensional vector space has $r$ elements, i.e., every basis has $r$ elements).
        Deduce Burnside's Basis Theorem: a set $y_1, \ldots, y_s$ is a minimal generating set for $P$ if and only if $\bar{y_1}, \ldots, \bar{y_s}$ is a basis of $\bar{P} = P/\Phi(P)$.
        Deduce that any minimal generating set for $P$ has $r$ elements.
        \item Prove that if $P/\Phi(P)$ is cyclic then $P$ is cyclic.
        Deduce that if $P/P'$ is cyclic then so is $P$.
        \item Let $\sigma$ be any automorphism of $P$ of prime order $q$ with $q \neq p$.
        Show that if $\sigma$ fixes the coset $x\Phi(P)$ then $\sigma$ fixes some element of this coset (note that since $\Phi(P)$ is characteristic in $P$ every automorphism of $P$ induces an automorphism of $P/\Phi(P)$).
        [Use the observation that $\sigma$ acts a permutation of order $1$ or $q$ on the $p^a$ elements in the coset $x\Phi(P)$.]
        \item Use parts (e) and (c) to deduce that every nontrivial automorphism of $P$ of order prime to $p$ induces a nontrivial automorphism on $P/\Phi(P)$.
        Deduce that any group of automorphisms of $P$ which has order prime to $p$ is isomorphic to a subgroup of $\mathrm{Aut}(\bar{P}) = GL_r(\fbb_p)$.
    \end{subexercises}
\end{exercise}

\begin{exercise}{6.1.27}
    Generalize part (d) of the preceding exercise as follows: let $p$ be a prime, let $P$ be a $p$-group and let $\bar{P} = P/\Phi(P)$ be elementary Abelian of order $p^r$.
    Prove that $P$ has exactly $\dfrac{p^r - 1}{p - 1}$ maximal subgroups.
    [Since every maximal subgroup of $P$ contains $\Phi(P)$, the maximal subgroups of $P$ are, by the Lattice Isomorphism Theorem, in bijective correspondence with the maximal subgroups of the elementary Abelian group $\bar{P}$.
    It therefore suffices to show that the number of maximal subgroups of an elementary Abelian $p$-group of order $p^r$ is as stated above.
    One way of doing this is to use the result that an Abelian group is isomorphic to its dual group (cf. Exercise 14 in Section 5.2) so the number of subgroups of index $p$ equals the number of subgroups of order $p$.]
\end{exercise}

\begin{exercise}{6.1.28}
    Prove that if $p$ is a prime and $P = Z_p \times Z_{p^2}$ then $|\Phi(P)| = p$ and $P/\Phi(P) \cong Z_p \times Z_p$.
    Deduce that $P$ has $p+1$ maximal subgroups.
\end{exercise}

\begin{exercise}{6.1.29}
    Prove that if $p$ is a prime and $P$ is a non-Abelian group of order $p^3$ then $\Phi(P) = Z(P)$ and $P/\Phi(P) \cong Z_p \times Z_p$.
    Deduce that $P$ has $p+1$ maximal subgroups.
\end{exercise}

\begin{exercise}{6.1.30}
    Let $p$ be an odd prime, let $P_1 = Z_p \times Z_{p^2}$ and let $P_2$ be the non-Abelian group of order $p^3$ which has an element of order $p^2$.
    Prove that $P_1$ and $P_2$ have the same lattice of subgroups.
\end{exercise}

\begin{exercise}{6.1.31}
    For any group $G$ a \emph{minimal normal subgroup} is a normal subgroup $M$ of $G$ such that the only normal subgroups of $G$ which are contained in $M$ are $1$ and $M$.
    Prove that every minimal normal subgroup of a finite solvable group is an elementary Abelian $p$-group for some prime $p$.
    [If $M$ is a minimal normal subgroup of $G$, consider its characteristic subgroups: $M'$ and $\langle x^p \mid x \in M \rangle$.]
\end{exercise}

\begin{exercise}{6.1.32}
    Prove that every maximal subgroup of a finite solvable group has prime power index.
    [Let $H$ be a maximal subgroup of $G$ and let $M$ be a minimal normal subgroup of $G$ — cf. the preceding exercise.
    Apply induction to $G/M$ and consider separately the two cases: $M \leq H$ and $M \not\leq H$.]
\end{exercise}

\begin{exercise}{6.1.33}
    Let $\pi$ be any set of primes.
    A subgroup $H$ of a finite group is called a \emph{Hall $\pi$-subgroup} of $G$ if the only primes dividing $|H|$ are in the set $\pi$ and $|H|$ is relatively prime to $|G:H|$.
    (Note that if $\pi = \{p\}$, Hall $\pi$-subgroups are the same as Sylow $p$-subgroups. Hall subgroups were introduced in Exercise 10 of Section 3.3).
    Prove the following generalization of Sylow's Theorem for solvable groups: if $G$ is a finite solvable group then for every set $\pi$ of primes, $G$ has a Hall $\pi$-subgroup and any two Hall $\pi$-subgroups (for the same set $\pi$) are conjugate in $G$.
    [Fix $\pi$ and proceed by induction on $|G|$, proving both existence and conjugacy at once.
    Let $M$ be a minimal normal subgroup of $G$, so $M$ is a $p$-group for some prime $p$.
    If $p \in \pi$, apply induction to $G/M$.
    If $p \notin \pi$, reduce to the case $|G| = p^\alpha n$, where $p^\alpha = |M|$ and $n$ is the order of a Hall $\pi$-subgroup of $G$.
    In this case let $N/M$ be a minimal normal subgroup of $G/M$, so $N/M$ is a $q$-group for some prime $q \neq p$.
    Let $Q \in \syl_q(N)$.
    If $Q \nsub G$ argue as before with $Q$ in place of $M$.
    If $Q$ is not normal in $G$, use Frattini's Argument to show $N_G(Q)$ is a Hall $\pi$-subgroup of $G$ and establish conjugacy in this case too.]
\end{exercise}

\begin{exercise}{6.1.34}
    Let $p$ be a prime dividing the order of the finite solvable group $G$.
    Assume $G$ has no nontrivial normal subgroups of order prime to $p$.
    Let $P$ be the largest normal $p$-subgroup of $G$ (cf. Exercise 37, Section 4.5).
    Note that Exercise 31 above shows that $P \neq 1$.
    Prove that $C_G(P) \leq P$, i.e., $C_G(P) = Z(P)$.
    [Let $N = C_G(P)$ and use the preceding exercise to show $N = Z(P) \times H$ for some Hall $\pi$-subgroup $H$ of $N$ — here $\pi$ is the set of all prime divisors of $|N|$ except for $p$.
    Show $H \nsub G$ to obtain the desired conclusion: $H = 1$.]
\end{exercise}

\begin{exercise}{6.1.35}
    Prove that if $G$ is a finite group in which every proper subgroup is nilpotent, then $G$ is solvable.
    [Show that a minimal counterexample is simple.
    Let $M$ and $N$ be distinct maximal subgroups chosen with $|M \cap N|$ as large as possible and apply Part 2 of Theorem 3 to show that $M \cap N = 1$.
    Now apply the methods of Exercise 53 in Section 4.5.]
\end{exercise}

\begin{exercise}{6.1.36}
    Let $p$ be a prime, let $V$ be a nonzero finite dimensional vector space over the field of $p$ elements and let $\varphi$ be an element of $GL(V)$ of order a power of $p$ (i.e., $V$ is a nontrivial elementary Abelian $p$-group and $\varphi$ is an automorphism of $V$ of $p$-power order).
    Prove that there is some nonzero element $v \in V$ such that $\varphi(v) = v$, i.e., $\varphi$ has a nonzero fixed point on $V$.
\end{exercise}

\begin{exercise}{6.1.37}
    Let $V$ be a finite dimensional vector space over the field of $2$ elements and let $\varphi$ be an element of $GL(V)$ of order $2$. (i.e., $V$ is a nontrivial elementary Abelian $2$-group and $\varphi$ is an automorphism of $V$ of order $2$).
    Prove that the map $v \mapsto v + \varphi(v)$ is a homomorphism from $V$ to itself.
    Show that every element in the image of this map is fixed by $\varphi$.
    Deduce that the subspace of elements of $V$ which are fixed by $\varphi$ has dimension $\geq \frac{1}{2}(\text{dimension } V)$.
    (Note that if $G$ is the semidirect product of $V$ with $\langle \varphi \rangle$, where $V \nsub G$ and $\varphi$ acts by conjugation on $V$ by sending each $v \in V$ to $\varphi(v)$, then the fixed points of $\varphi$ on $V$ are $C_V(\varphi)$ and the above map is simply the commutator map: $v \mapsto [v, \varphi]$.
    In this terminology the problem is to show that $|C_V(\varphi)|^2 \geq |V|$.)
\end{exercise}

\begin{exercise}{6.1.38}
    Use the preceding exercise to prove that if $P$ is a $2$-group which has a cyclic center and $M$ is a subgroup of index $2$ in $P$, then the center of $M$ has rank $\leq 2$.
    [The group $G/M$ of order $2$ acts by conjugation on the $\fbb_2$ vector space: $\{z \in Z(M) \mid z^2 = 1\}$ and the fixed points of this action are in the center of $P$.]
\end{exercise}